\chapter{Tachypnea}

\section*{Definition}
Tachypnea is an abnormally rapid respiratory rate, defined in adults as a rate greater than 20 breaths/min at rest, and in children by age-specific thresholds (greater than 60/min in neonates, greater than 40/min in children 1--5 years). It is a sensitive, non-specific sign of underlying physiologic stress.

\section*{Pathophysiology}
Tachypnea results from activation of central or peripheral chemoreceptors (in response to hypoxemia, hypercapnia, or acidosis), pulmonary stretch receptors, or psychogenic stimulation. It may be compensatory (maintaining minute ventilation when tidal volume is reduced) or pathologic (wasted breathing).

\section*{Etiology}
\begin{itemize}
  \item \textbf{Pulmonary:} Pneumonia, asthma, COPD exacerbation, pulmonary embolism, pneumothorax, ARDS, interstitial lung disease
  \item \textbf{Cardiac:} Heart failure (pulmonary edema), pericardial tamponade, congenital heart disease
  \item \textbf{Metabolic:} Diabetic ketoacidosis (Kussmaul breathing), lactic acidosis, salicylate poisoning, renal failure
  \item \textbf{Fever/Infection:} Any febrile illness, sepsis (early compensatory phase)
  \item \textbf{Psychogenic:} Panic attack, anxiety disorder, hyperventilation syndrome
  \item \textbf{Other:} Anemia, hyperthyroidism, pain, drug withdrawal, altitude sickness
\end{itemize}

\section*{Indications for Evaluation}
Seek care if:
\begin{itemize}
  \item Tachypnea with oxygen saturation less than 92\% or signs of respiratory distress
  \item Rapid onset associated with chest pain, hemoptysis, or unilateral leg swelling (suggests PE)
  \item Altered mental status, confusion, or lethargy
  \item Fever with toxic appearance or signs of sepsis
  \item Tachypnea with fruity breath odor (DKA) or known diabetes
\end{itemize}

\section*{Related Symptoms}
\begin{itemize}
  \item Dyspnea, shallow breathing, accessory muscle use, tachycardia, palpitations, chest tightness, anxiety
\end{itemize}
\textbf{Note:} In children, tachypnea is often the earliest and most reliable sign of pneumonia, routinely preceding auscultatory findings.

\section*{Self-Care / Home Management}
\begin{itemize}
  \item Measure respiratory rate accurately (count for 30 seconds, multiply by 2) in a resting state
  \item Sit upright and practice slow, controlled breathing (4 seconds in, 6 seconds out)
  \item Ensure adequate hydration; avoid stimulants (caffeine) that may exacerbate tachypnea
  \item If hyperventilation suspected, low-normal CO2 is preferable to rebreathing into a paper bag (risk of hypoxia)
  \item Monitor oxygen saturation with home pulse oximeter if available; seek care if persistent or worsening
\end{itemize}

\section*{Diagnostic Approach}
\begin{itemize}
  \item Complete history (onset, duration, triggers, associated symptoms) and full vital signs
  \item Arterial blood gas or venous blood gas to assess pH, pCO2, pO2, and bicarbonate
  \item Chest X-ray for pulmonary pathology
  \item ECG for cardiac ischemia, arrhythmia, or pulmonary embolism signs (S1Q3T3)
  \item Laboratory: CBC, basic metabolic panel, troponin, D-dimer, thyroid function tests as indicated
\end{itemize}

\section*{Treatment / Management Overview}
\begin{itemize}
  \item Treat the underlying cause rather than suppressing the respiratory drive directly
  \item Supplemental oxygen to correct hypoxemia if present (target SpO2 greater than 92\%)
  \item Bronchodilators for reactive airway disease; diuretics for pulmonary edema
  \item Insulin and fluid resuscitation for DKA; sodium bicarbonate rarely indicated for acidosis
  \item Antipyretics for fever-related tachypnea; benzodiazepines cautiously for psychogenic hyperventilation
\end{itemize}

\clearpage
